% Part: lambda-calculus
% Chapter: introduction
% Section: church-rosser

\documentclass[../../../include/open-logic-section]{subfiles}

\begin{document}

\olfileid{lam}{int}{cr}
\olsection{خاصية تشيرش--روسر}

\begin{thm}
\ollabel{thm:church-rosser}
لتكن~$M$ و$N_1$ و$N_2$ حدودًا بحيث~$M \red N_1$ و$M \red N_2$.
عندئذ يوجد حد~$P$ بحيث~$N_1 \red P$ و$N_2 \red P$.
\end{thm}

\begin{cor}
افترض أن~$M$ يمكن اختزاله إلى صورة سوية. عندئذ تكون هذه الصورة السوية
وحيدة.
\end{cor}

\begin{proof}
إذا كان~$M \red N_1$ و$M \red N_2$، فبالمبرهنة السابقة يوجد حد~$P$
يختزل كل من~$N_1$ و$N_2$ إلى~$P$. وإذا كان كل من~$N_1$ و$N_2$ في صورة
سوية، فلا يمكن أن يحدث هذا إلا إذا~$N_1 \ident P \ident N_2$.
\end{proof}

وأخيرًا، سنقول إن الحدين~$M$ و$N$ \emph{متكافئان بحسب~$\beta$}، أو
\emph{متكافئان} فحسب، إذا كانا يختزلان إلى حد مشترك؛ أي إذا وجد~$P$ ما
بحيث~$M \red P$ و$N \red P$. ونكتب هذا~$M \equal[\beta] N$. ويمكنك
باستعمال~\olref{thm:church-rosser} التحقق من أن~$\equal[\beta]$ علاقة
تكافؤ لها الخاصية الإضافية الآتية: لكل~$M$ و$N$، إذا كان~$M \red N$
أو~$N \red M$، فإن~$M \equal[\beta] N$. (بل يمكن إثبات أن
$\equal[\beta]$ هي \emph{أصغر} علاقة تكافؤ لها هذه الخاصية.)

\end{document}
